\section{Method}

A semantic cache stores a set of (query, answer) pairs $C$. When a new query $q$ arrives, the cache retrieves the $k$ nearest stored entries by cosine similarity in the embedding space. The core decision is whether to reuse the response associated with the nearest entry $c_{(1)}$ or to forward $q$ to the LLM. This is a binary choice: serve from cache ($\hat{y}=1$) or defer to the model ($\hat{y}=0$). The ground-truth label $y=1$ indicates that $c_{(1)}$ is a genuine semantic duplicate of $q$ and that serving its cached response is correct; $y=0$ means the query is distinct and serving the nearest entry's answer would be an error.

SERVE makes this reuse decision using four features, all of which are byproducts of the top-$k$ nearest-neighbor search that every semantic cache already performs. Let $s_1 = \mathrm{sim}(q, c_{(1)})$ denote the cosine similarity to the nearest cache entry and $s_2 = \mathrm{sim}(q, c_{(2)})$ the similarity to the runner-up. The margin $m = s_1 - s_2$ measures how contested the best match is: a large $m$ means the top match stands clearly apart, while $m$ near zero signals ambiguity and a genuine risk that a different, semantically distinct entry is nearly as close. The fourth feature is the local density $\rho$, defined as the mean similarity over the $k=8$ nearest cache entries.
% E1
Together these four quantities capture both the absolute quality of the best match and the geometry of the query's local neighborhood. They expose structural information about whether $s_1$ reflects a clean, isolated match or a crowded region where cosine proximity is ambiguous. Neither the runner-up similarity, the margin, nor the local density is visible from $s_1$ alone, yet all three are computed and then discarded by the top-$k$ search that every deployed semantic cache already runs.

These four quantities contain an exact linear dependence ($m = s_1 - s_2$), so the feature space carries only three independent degrees of freedom. For a linear model this would yield infinite variance inflation factors on $s_1$, $s_2$, and $m$, rendering coefficient estimates unstable. Gradient-boosted trees, however, split on axis-aligned thresholds rather than inverting covariance matrices; the linear dependence therefore does not inflate variance or corrupt the split-search procedure. The model remains free to use whichever of the four representations best separates the data at a given node, and the redundancy exposes the structural relationships (absolute similarity, margin, and runner-up quality) directly rather than requiring the tree to compose them through successive splits.

The gate is a \texttt{GradientBoostingClassifier} \cite{friedman2001greedy} with 150 estimators and maximum tree depth 3, trained to predict the binary label $y$ from the four-dimensional feature vector $\mathbf{x} = [s_1, m, s_2, \rho]$.
% E1
During inference, the classifier outputs a score $g_\theta(\mathbf{x}) \in [0,1]$, which is compared against a decision threshold $\tau_g$. The gate serves from cache when $g_\theta(\mathbf{x}) > \tau_g$ and defers to the LLM otherwise. The threshold $\tau_g$ is not learned during training; it is swept post-hoc on a calibration set to select the operating point that satisfies a target false-serve budget, giving the deployer direct control over the tradeoff between hit rate and reliability.

The classifier is trained on a labeled set of (feature vector, label) pairs extracted from a held-out gate-training split that is disjoint from both the cache population and the final test set. For each query in this split, the retrieval geometry against the full cache is computed, yielding $\mathbf{x}$, and the label $y$ is determined from human-annotated duplicate-pair ground truth. This protocol matches the deployment scenario: the gate learns to predict correctness from retrieval features alone, without access to the text content of queries or answers. Training minimizes log-loss as a standard binary classification objective. The operating threshold $\tau_g$ is then calibrated separately to satisfy the false-serve budget chosen by the deployer.

The computational cost of the gate is negligible. Because every feature in $\mathbf{x}$ is already computed during the top-$k$ nearest-neighbor search that the cache performs regardless of the gating mechanism, SERVE adds no additional embedding calls or index lookups. The only supplementary cost is evaluating the 150 tree ensemble, which takes 0.8\,$\mu$s per query on an Apple M-series CPU.
% E5
At this latency the gate is effectively free relative to the cost of embedding computation, nearest-neighbor search, and LLM inference.

A structural consequence of this design is that the gate is fully decoupled from the cache backend. SERVE does not modify the cache's data structure, eviction policy, or embedding model. It consumes retrieval features that any nearest-neighbor cache produces and emits a binary serve/defer decision. The gate can be dropped into GPTCache, MeanCache, vCache, or any custom semantic cache without altering the underlying system, and it can be retrained on whatever embedding model and dataset the cache uses. The false-serve budget that controls $\tau_g$ is the only deployment-time knob. The following section describes the experimental protocol used to evaluate the gate against fixed-threshold and learned baselines under this budget-controlled regime.
